% Options for packages loaded elsewhere
\PassOptionsToPackage{unicode}{hyperref}
\PassOptionsToPackage{hyphens}{url}
\documentclass[
  ignorenonframetext,
]{beamer}
\newif\ifbibliography
\usepackage{pgfpages}
\setbeamertemplate{caption}[numbered]
\setbeamertemplate{caption label separator}{: }
\setbeamercolor{caption name}{fg=normal text.fg}
\beamertemplatenavigationsymbolsempty
% remove section numbering
\setbeamertemplate{part page}{
  \centering
  \begin{beamercolorbox}[sep=16pt,center]{part title}
    \usebeamerfont{part title}\insertpart\par
  \end{beamercolorbox}
}
\setbeamertemplate{section page}{
  \centering
  \begin{beamercolorbox}[sep=12pt,center]{section title}
    \usebeamerfont{section title}\insertsection\par
  \end{beamercolorbox}
}
\setbeamertemplate{subsection page}{
  \centering
  \begin{beamercolorbox}[sep=8pt,center]{subsection title}
    \usebeamerfont{subsection title}\insertsubsection\par
  \end{beamercolorbox}
}
% Prevent slide breaks in the middle of a paragraph
\widowpenalties 1 10000
\raggedbottom
\AtBeginPart{
  \frame{\partpage}
}
\AtBeginSection{
  \ifbibliography
  \else
    \frame{\sectionpage}
  \fi
}
\AtBeginSubsection{
  \frame{\subsectionpage}
}
\usepackage{iftex}
\ifPDFTeX
  \usepackage[T1]{fontenc}
  \usepackage[utf8]{inputenc}
  \usepackage{textcomp} % provide euro and other symbols
\else % if luatex or xetex
  \usepackage{unicode-math} % this also loads fontspec
  \defaultfontfeatures{Scale=MatchLowercase}
  \defaultfontfeatures[\rmfamily]{Ligatures=TeX,Scale=1}
\fi
\usepackage{lmodern}
\ifPDFTeX\else
  % xetex/luatex font selection
\fi
% Use upquote if available, for straight quotes in verbatim environments
\IfFileExists{upquote.sty}{\usepackage{upquote}}{}
\IfFileExists{microtype.sty}{% use microtype if available
  \usepackage[]{microtype}
  \UseMicrotypeSet[protrusion]{basicmath} % disable protrusion for tt fonts
}{}
\makeatletter
\@ifundefined{KOMAClassName}{% if non-KOMA class
  \IfFileExists{parskip.sty}{%
    \usepackage{parskip}
  }{% else
    \setlength{\parindent}{0pt}
    \setlength{\parskip}{6pt plus 2pt minus 1pt}}
}{% if KOMA class
  \KOMAoptions{parskip=half}}
\makeatother
\usepackage{longtable,booktabs,array}
\newcounter{none} % for unnumbered tables
\usepackage{calc} % for calculating minipage widths
\usepackage{caption}
% Make caption package work with longtable
\makeatletter
\def\fnum@table{\tablename~\thetable}
\makeatother
\setlength{\emergencystretch}{3em} % prevent overfull lines
\providecommand{\tightlist}{%
  \setlength{\itemsep}{0pt}\setlength{\parskip}{0pt}}
\usepackage{bookmark}
\IfFileExists{xurl.sty}{\usepackage{xurl}}{} % add URL line breaks if available
\urlstyle{same}
\hypersetup{
  hidelinks,
  pdfcreator={LaTeX via pandoc}}

\author{\texorpdfstring{}{}}
\date{}

\begin{document}

\section{Program analysis for machine learning and comparison
tables}\label{program-analysis-for-machine-learning-and-comparison-tables}

\begin{frame}[fragile]{Program analysis for ML}
\protect\phantomsection\label{program-analysis-for-ml}
Several systems address the intersection of program analysis and machine
learning that COGANT operates in:

\textbf{code2seq and code2vec} {[}@alon2019code2vec{]} represent
programs as sets of AST paths between terminal nodes. These path-based
representations are effective for method naming and code summarization
but discard the full graph topology that graph-neural-network-based
models exploit. COGANT preserves the complete program graph ---
including control-flow and data-flow edges --- and can export it in
tensor form, enabling downstream architectures that reason over richer
relational structure.

\textbf{Learning to Represent Programs with Graphs}
{[}@allamanis2018learning{]} is the closest conceptual neighbour: it
constructs typed graphs fusing AST, control, and data-flow edges and
feeds them to a gated graph neural network for variable-misuse and
variable-naming tasks. COGANT's program graph IR generalises that
construction into a pluggable and confidence-annotated pipeline, and its
export contract is designed so that the tensor output remains directly
compatible with gated-message-passing architectures in the same family.

\textbf{Gated Graph Neural Networks (GGNN)} {[}@li2016gated{]}
introduced gated recurrent propagation for graph-structured program
representations, demonstrating strong results on variable misuse
detection and other code tasks. COGANT's export contract (PyG
\texttt{Data} objects with typed \texttt{edge\_index} and
\texttt{edge\_attr} {[}@fey2019pyg{]}) is directly compatible with
GGNN-style models; the kind and role indices provide the discrete node
and edge types that typed message-passing layers require.

\textbf{Typilus} {[}@allamanis2020typilus{]} and \textbf{LambdaNet}
{[}@wei2020lambdanet{]} use graph neural networks to predict type
annotations from structural program graphs. Both consume exactly the
kind of typed adjacency structure COGANT emits, and both rely on message
passing over node kinds similar to COGANT's \texttt{NodeKind} taxonomy,
so COGANT's exports can serve as an upstream graph generator for
Typilus- and LambdaNet-style type inference.

\textbf{CodeQL} (Semmle/GitHub), whose declarative query language QL is
formalised in {[}@avgustinov2016ql{]}, provides a declarative query
language over relational representations of code. While CodeQL excels at
security analysis with hand-written queries, its outputs are query
results rather than tensor-ready graph bundles. COGANT occupies the
complementary niche: it produces the graph data that learned models
consume, and could ingest CodeQL query results as an additional evidence
source feeding the confidence model.

\textbf{CodeBERT} {[}@feng2020codebert{]} and related pre-trained models
operate at the token level, learning representations from natural
language and code jointly. \textbf{GraphCodeBERT}
{[}@guo2021graphcodebert{]} extends this line by injecting data-flow
edges into the pre-training objective, which shows that even token-level
models benefit from the kinds of data-flow relationships COGANT surfaces
first-class in its program graph. These models are complementary to
COGANT's graph-centric approach: their embeddings can serve as optional
node features in COGANT's Generalized Notation Notation (GNN) export
(the export schema already reserves dimensions for text embeddings as
documented in \texttt{../cogant/docs/export/README.md}).

\begin{block}{Feature matrix: COGANT vs.~related tools}
\protect\phantomsection\label{feature-matrix-cogant-vs.-related-tools}
The following matrix contrasts COGANT's capabilities with the related
tools discussed in this section. Entries marked ``yes'' indicate
first-class support; ``partial'' indicates limited or indirect support;
``no'' indicates the feature is out of scope for that tool.

\textbf{Table 8. Feature comparison of program-to-model toolchains.}

{\def\LTcaptype{none} % do not increment counter
\begin{longtable}[]{@{}
  >{\raggedright\arraybackslash}p{(\linewidth - 10\tabcolsep) * \real{0.1765}}
  >{\centering\arraybackslash}p{(\linewidth - 10\tabcolsep) * \real{0.1569}}
  >{\centering\arraybackslash}p{(\linewidth - 10\tabcolsep) * \real{0.1961}}
  >{\centering\arraybackslash}p{(\linewidth - 10\tabcolsep) * \real{0.1176}}
  >{\centering\arraybackslash}p{(\linewidth - 10\tabcolsep) * \real{0.1569}}
  >{\centering\arraybackslash}p{(\linewidth - 10\tabcolsep) * \real{0.1961}}@{}}
\toprule\noalign{}
\begin{minipage}[b]{\linewidth}\raggedright
Feature
\end{minipage} & \begin{minipage}[b]{\linewidth}\centering
COGANT
\end{minipage} & \begin{minipage}[b]{\linewidth}\centering
code2vec
\end{minipage} & \begin{minipage}[b]{\linewidth}\centering
GGNN
\end{minipage} & \begin{minipage}[b]{\linewidth}\centering
CodeQL
\end{minipage} & \begin{minipage}[b]{\linewidth}\centering
CodeBERT
\end{minipage} \\
\midrule\noalign{}
\endhead
Full program graph (AST + CFG + DFG) & yes & no & input-only & yes &
no \\
Typed node/edge taxonomy & yes & no & partial & yes & no \\
Confidence scoring per assertion & yes & no & no & no & no \\
Provenance tracking & yes & no & no & partial & no \\
State-space extraction & yes & no & no & no & no \\
Temporal regime classification & yes & no & no & no & no \\
Dynamic enrichment (coverage, traces) & yes & no & no & partial & no \\
Generalized Notation Notation output & yes & no & no & no & no \\
Tensor export (PyG, DGL, HDF5) & yes & partial & input-only & no & no \\
Pluggable translation rules & yes & no & no & yes & no \\
Human review loop & yes & no & no & partial & no \\
Multi-language front-ends & partial (Python; JS/TS optional) & yes & no
& yes & yes \\
\bottomrule\noalign{}
\end{longtable}
}

COGANT is distinct from the other toolchains in three ways: first, it
explicitly models uncertainty through confidence tiers tied to evidence
provenance; second, it produces a structured Active Inference notation
as its primary output rather than an opaque tensor; and third, it
composes static and dynamic evidence in a single pipeline rather than
specializing to one.
\end{block}

\begin{block}{Input/output comparison vs prior art}
\protect\phantomsection\label{inputoutput-comparison-vs-prior-art}
Table 8 contrasts fine-grained feature flags; Table 9 expands the frame
to include the \emph{input/output contract} of each approach, because
the most consequential difference between COGANT and its neighbours is
what a user has to supply (training data, hand-written queries, manual
modelling) and what they get back (vector, query table, simulator-ready
model). The comparison covers code-representation learning (code2vec),
learned graph models for programs (GGNN, Typilus, LambdaNet),
code-property-graph-based analysers (CodeQL, the original Joern/CPG
line), compiler IRs (PDG, LLVM IR, MLIR), and Active Inference tooling
(hand-authored GNN with PyMDP as the downstream runtime).

\textbf{Table 9. Input/output comparison of COGANT and prior
approaches.}

{\def\LTcaptype{none} % do not increment counter
\begin{longtable}[]{@{}
  >{\raggedright\arraybackslash}p{(\linewidth - 10\tabcolsep) * \real{0.1364}}
  >{\raggedright\arraybackslash}p{(\linewidth - 10\tabcolsep) * \real{0.1364}}
  >{\raggedright\arraybackslash}p{(\linewidth - 10\tabcolsep) * \real{0.1364}}
  >{\centering\arraybackslash}p{(\linewidth - 10\tabcolsep) * \real{0.2273}}
  >{\raggedright\arraybackslash}p{(\linewidth - 10\tabcolsep) * \real{0.1364}}
  >{\centering\arraybackslash}p{(\linewidth - 10\tabcolsep) * \real{0.2273}}@{}}
\toprule\noalign{}
\begin{minipage}[b]{\linewidth}\raggedright
Approach
\end{minipage} & \begin{minipage}[b]{\linewidth}\raggedright
Primary input
\end{minipage} & \begin{minipage}[b]{\linewidth}\raggedright
Primary output
\end{minipage} & \begin{minipage}[b]{\linewidth}\centering
Requires training
\end{minipage} & \begin{minipage}[b]{\linewidth}\raggedright
Languages (as shipped)
\end{minipage} & \begin{minipage}[b]{\linewidth}\centering
Produces Active Inference model
\end{minipage} \\
\midrule\noalign{}
\endhead
\textbf{COGANT} (this work) & Source repository (checkout or URL) &
Generalized Notation Notation bundle (A/B/C/D, state space, Markov
blanket, tensor views) & no & Python; JS/TS optional
(\texttt{cogant{[}multilang{]}} + grammars) & yes (end-to-end) \\
code2vec / code2seq {[}@alon2019code2vec{]} & Single method or function
body & Fixed-size embedding vector (predicted method name or tag) & yes
(14M-method corpus) & Java (primary), C\#, Python (partial) & no \\
Gated GNN for programs {[}@allamanis2018learning; @li2016gated{]} &
Typed program graph (AST + control + data-flow) & Task-specific
prediction (variable misuse, variable naming) & yes (task-specific
labels) & C\# (original), Java & no \\
Typilus {[}@allamanis2020typilus{]} / LambdaNet {[}@wei2020lambdanet{]}
& Typed program graph & Predicted type annotations & yes & Python /
TypeScript & no \\
Program Dependence Graph {[}ferrante1987pdg; horwitz1990slicing{]} &
Single procedure or interprocedural bundle & PDG / System Dependence
Graph & no & Any (formalism-level) & no \\
LLVM / MLIR IR {[}@lattner2004llvm; @lattner2021mlir{]} & Source in
supported front-end language & SSA-form compiler IR, optimisation
passes, code generation & no & C/C++/Rust/Swift/many via LLVM & no \\
CodeQL / QL {[}@avgustinov2016ql{]} & Source repository + hand-written
query & Query result table (alerts, findings) & no & Python, JS/TS,
Java, C\#, Go, C/C++ & no \\
CodeBERT / GraphCodeBERT {[}@feng2020codebert; @guo2021graphcodebert{]}
& Token (and DFG) sequence for a code fragment & Contextual embeddings
for downstream tasks & yes (multi-million-pair corpus) & Python, Java,
JS, PHP, Ruby, Go & no \\
PyMDP {[}@heins2022pymdp{]} & Hand-authored A/B/C/D matrices (Python
objects) & Active Inference simulation trajectories & no & N/A (runtime,
not extractor) & yes (consumer of hand-authored input) \\
Generalized Notation Notation reference {[}@friedman2024gnn{]} &
Hand-authored GNN Markdown or JSON & State-space/process model artifacts
& no & N/A (notation + validator) & yes (format, not extractor) \\
\bottomrule\noalign{}
\end{longtable}
}

Three things are visible in this table that the fine-grained feature
matrix does not capture. First, \textbf{COGANT is the only row whose
input is a raw repository and whose output is a simulator-ready Active
Inference model}: every other Active-Inference entry in the rightmost
column (PyMDP, the GNN reference) requires a human to author the model
by hand, and every code-modelling entry (code2vec through CodeBERT)
produces either a vector, a type annotation, or a query result rather
than a generative model. Second, \textbf{COGANT's rule-based pipeline
does not require training}, which places it alongside the compiler-IR
and code-property-graph lines rather than the learned-embedding lines in
Section 8's ``training'' column. Third, \textbf{the languages column
highlights that COGANT's v0.5.x front-end set (Python first-class;
JavaScript / TypeScript via optional \texttt{cogant{[}multilang{]}} and
\texttt{tree-sitter} when installed) is a deliberate scope choice, not a
structural limitation}: the rule engine and state-space compiler consume
a language-agnostic \texttt{ProgramGraph} IR, so adding a further parser
(Java, Go, Rust, C/C++) is a matter of implementing the plugin interface
in \texttt{../cogant/docs/plugins/README.md} and does not touch the
translation, matrix, or export layers. The
\texttt{examples/zoo/13\_js\_observer} cross-language round-trip (§5 and
\texttt{cogant/docs/evaluation/ROUNDTRIP\_IMPROVEMENT.md}) establishes
the template for validating new parsers before release.
\end{block}
\end{frame}

\end{document}
